\section{What \odag{} can learn from \fca{}}
\label{sec:learn}

\fca{} has forty years of theory; \odag{} has a working system and a short
list of invariants it will not give up. This section goes through the parts of
the theory that could be imported, and for each asks four questions:

\begin{enumerate}[label=(\alph*),leftmargin=2.0em,itemsep=1pt]
\item what is it, in one paragraph a first-year reader can follow;
\item what would \odag{} gain;
\item what would it cost --- specifically, which of \odag{}'s guarantees
  (\S\ref{sec:guarantees}) does it threaten;
\item verdict.
\end{enumerate}

The verdicts use a fixed vocabulary, because the distinction between them is
the most transferable thing in this section:

\begin{center}\small
\begin{tabularx}{0.94\textwidth}{@{}lX@{}}
\toprule
\textbf{Adopt (core)} & belongs in the shared, hashed, merged graph\\
\textbf{Adopt (derived)} & belongs in a separate local index: regenerable,
  evictable, never merged, outside the canonical root\\
\textbf{Adopt (surface)} & belongs in the interface layer --- rendering,
  suggestion, tooling --- and touches no stored bytes\\
\textbf{Research} & desirable, but breaks something until a real problem is
  solved\\
\textbf{Refuse} & incompatible with an invariant \odag{} exists to provide\\
\bottomrule
\end{tabularx}
\end{center}

That vocabulary is itself a lesson learned from \fca{}: the reason \odag{} can
consider adopting so much is that it has a place to put derived things where
they cannot contaminate the canonical form.

\subsection{Implications and the stem base}
\label{sec:learn-implications}

\paragraph{(a) What it is.} An implication $B \to C$ holds when everything
having all of $B$ also has all of $C$ (\S\ref{sec:implications}). The
Duquenne--Guigues \emph{stem base} \citep{guigues1986} is the unique smallest
set of implications generating all the others. It can be computed by
NextClosure at the same time as the concepts \citep{ganter1984}.

\paragraph{(b) What \odag{} would gain.} Two genuinely useful things.

First, a \textbf{suggestion engine} for classification --- which is exactly
the gap \odag{}'s contract admits it has (``classification is the higher
layer's job''). Measured on the running example, \fca{} reports
$\textsf{Refundable} \to \textsf{Booked}, \textsf{Japan}$. If a user now files
a refundable document \emph{without} marking it \textsf{Booked}, the system
can say: \emph{every other refundable thing here is booked --- did you mean to
mark this one?} That is a well-founded prompt, computed from the store's own
contents, needing no language model and no external ontology.

Second, a \textbf{redundancy lint}. If $\textsf{Hotel} \to \textsf{Travel}$
holds extensionally for every hotel in the store, someone may have forgotten
to assert the edge $\textsf{Hotel} \to \textsf{Travel}$ in the graph. The
implication is evidence for a missing subsumption.

\paragraph{(c) What it would cost.} Storing implications in the shared graph
would be a category error, and it is worth being precise about why, because
the same argument recurs throughout this section.

An extensional implication is a \textbf{fact about today's population, not a
law}. $\textsf{Flight} \to \textsf{Booked}$ holds in the running example only
because nobody has yet filed an unbooked flight. Now consider what happens
under merge: a collaborator who has one unbooked flight merges their store
into yours, and the implication is destroyed. An implication is therefore
\emph{non-monotone} --- it can be falsified by the arrival of new information
--- which is exactly the property that guarantee G2 forbids in the shared
store. Stored as data, implications would make two replicas disagree, and
would make the canonical root depend on which subset of the world a replica
had seen.

There is a second, subtler cost: computing the stem base is expensive
(equivalent in the worst case to enumerating pseudo-intents, which is not
polynomial-delay), so it is a batch job, not something to run on every write.

\paragraph{(d) Verdict: \textbf{Adopt (derived)} for the lint, and
\textbf{Adopt (surface)} for the suggestion.} Compute implications in a
separate, local, regenerable index --- the same discipline \odag{} already
uses for cone summaries --- and surface them as \emph{prompts to a human or an
agent}, never as stored claims. If the user accepts a prompt, what gets stored
is an ordinary asserted edge with ordinary provenance. The implication was
evidence; the edge is the claim. Keeping that distinction is the difference
between a helpful system and one that silently invents knowledge.

\begin{keyidea}
The general rule this instance establishes: \textbf{\fca{}'s derived output
may enter \odag{} only as a suggestion to an author, never as an assertion by
the system.} Everything that follows in this section obeys it.
\end{keyidea}

\subsection{Conceptual scaling, and the levels of measurement}
\label{sec:learn-scaling}

\paragraph{(a) What it is.} Conceptual scaling (\S\ref{sec:fca-extensions})
turns non-binary data into a context by choosing, per attribute, which
comparisons matter: nominal (distinct values), ordinal (thresholds),
interordinal (both ends), and so on \citep[ch.~1]{ganter1999}.

\paragraph{(b) What \odag{} would gain.} Not a mechanism --- \odag{} already
has one --- but a \textbf{vocabulary and a completeness check}. \odag{}'s
dimension kinds turn out to line up exactly with a classification that
predates both systems: Stevens's levels of measurement \citep{stevens1946}.
Canonicalisation, in each case, is quotienting by the transformations the
scale considers meaningless:

\begin{center}\small
\begin{tabularx}{\textwidth}{@{}p{2.3cm} p{4.3cm} p{3.1cm} X@{}}
\toprule
\textbf{Stevens level} & \textbf{Meaningful transformations} &
\textbf{\fca{} scale} & \textbf{\odag{} kind}\\
\midrule
nominal  & any relabelling            & nominal      & plain nodes\\
ordinal  & any order-preserving map   & ordinal      & \emph{missing}\\
interval & affine $x \mapsto ax + b$  & interordinal & affine (temperatures)\\
ratio    & scaling $x \mapsto ax$     & interordinal & linear (weight, length)\\
absolute & none                       & nominal      & count\\
\bottomrule
\end{tabularx}
\end{center}

The table has one conspicuous hole, and finding holes is what a good
vocabulary is for: \odag{} has \emph{no ordinal kind}. There is no way to say
that \textsf{beginner} $<$ \textsf{intermediate} $<$ \textsf{expert} as a
declared dimension, so such chains have to be built as ordinary edges, which
works but is unindexed and gives no arithmetic. \fca{}'s ordinal scale is the
template for what the declaration would look like.

\paragraph{(c) What it would cost.} Adding an ordinal kind is cheap on both
admissibility axes: it is monotone, and the order is computable from the
declared chain plus the value names, so canonicality is preserved. The real
risk is social rather than technical --- an ordinal kind invites users to fake
arithmetic on ranks (``the average of beginner and expert''), which is exactly
the abuse Stevens wrote his paper to warn against.

\paragraph{(d) Verdict: \textbf{Adopt (core)}, carefully.} The scale
vocabulary should be documented as \odag{}'s own; the ordinal kind should be
added, with the arithmetic deliberately absent.

\subsection{Pattern structures: the theory \odag{} already uses}
\label{sec:learn-patterns}

\paragraph{(a) What it is.} \citet{ganter2001} generalise \fca{} so that an
object's description need not be a set of attributes. Instead descriptions
live in a meet-semilattice $(D, \sqcap)$ with a similarity operation. The
derivation operators become
\[
  A^{\square} = \bigsqcap_{g \in A} \delta(g),
  \qquad
  d^{\square} = \{g \in G : d \sqsubseteq \delta(g)\},
\]
and the whole theory --- concepts, lattice, implications --- goes through
unchanged. The most-used instance is the \emph{interval pattern structure}:
descriptions are intervals, $\sqcap$ is the convex hull, and $d \sqsubseteq e$
means $e \subseteq d$ as intervals \citep{kaytoue2011}.

\paragraph{(b) What \odag{} would gain.} Recognition, and then results.
\odag{}'s parametric dimensions (\S\ref{sec:dimensions}) \emph{are} an
interval pattern structure: \verb|weight(..5kg)| is an interval,
containment of denoted value sets is $\sqsubseteq$, and the star of values
under a dimension head is the set of descriptions in play. \odag{} arrived at
this independently, from the practical need to answer a courier's
``anything under five kilos?'' without pre-quantising.

Knowing that it is an instance of a studied formalism buys three things:

\begin{itemize}[leftmargin=1.4em]
\item \textbf{Correctness assurance.} The pattern-structure framework tells
  you exactly what conditions $\sqcap$ must satisfy for the theory to hold
  (associativity, commutativity, idempotence --- a semilattice). \odag{}'s
  requirement that only exact-arithmetic kinds may enter the canonical order
  is the same condition, discovered the hard way.
\item \textbf{Projections.} Pattern structures come with a theory of
  \emph{projections}: simplifying descriptions (coarsening intervals) in a way
  that provably yields a sound approximation of the concept lattice. This is
  the principled version of the ``derived index for hot dimensions'' idea that
  \odag{} currently treats as an engineering hunch.
\item \textbf{Multi-dimensional descriptions.} \citet{kaytoue2011} handle a
  \emph{vector} of intervals as one description --- which is precisely
  \odag{}'s open question about how several coordinates
  (\verb|from|, \verb|to|, \verb|duration|) should combine. The
  pattern-structure answer is the product semilattice with componentwise
  $\sqcap$, and it comes with the meet already defined.
\end{itemize}

\paragraph{(c) What it would cost.} Nothing, if imported as \emph{theory}.
Importing the \emph{machinery} --- computing the pattern concept lattice ---
would reintroduce the closure problem of Section~\ref{sec:not-closed}, with
the added hazard that interval pattern lattices are typically larger than
their binary counterparts.

\paragraph{(d) Verdict: \textbf{Adopt (theory)}, immediately and at no cost.}
Re-describe \odag{}'s dimensions in pattern-structure language in the design
documents; use the projection theory when the per-dimension index is finally
built; adopt the product-semilattice answer for multi-coordinate descriptions
when that question is decided. This is the single highest-value, lowest-cost
item on the list: it is free, and it converts an engineering hunch into a
known-correct construction.

\subsection{Stability, support, and knowing which concepts are worth having}
\label{sec:learn-stability}

\paragraph{(a) What it is.} Two ways of ranking concepts. \emph{Support} is
extent size --- keep the concepts with many members, which yields the iceberg
lattice \citep{stumme2002}. \emph{Stability} \citep{kuznetsov2007} is the
fraction of sub-collections of the extent that would still produce the same
intent --- keep the concepts that do not depend on any one or two objects.

\paragraph{(b) What \odag{} would gain.} Half of this machinery already
exists. The stored \verb|descendant_count| \emph{is} the extent size of a node, kept
current by delta on every write. So half of this machinery exists and is
already load-bearing (the query planner uses it to order cones). What is
missing is the \emph{use} of these numbers to decide what deserves to exist.

Concretely, three consumers:

\begin{enumerate}[leftmargin=1.7em]
\item \textbf{Which unnamed meets to suggest.} Of the four meets in
  Figure~\ref{fig:completion}, which should the system offer to name? Support
  and stability rank them. From Table~\ref{tab:concepts}, ``booked Japan
  documents'' has support 4 but stability $0.38$; ``refundable Japan
  bookings'' has lower support (3) but higher stability ($0.62$). The second
  is the better candidate for a category, and only stability says so.
\item \textbf{Which cones to index.} \odag{}'s published cone summaries
  currently use a fixed threshold rule on \verb|descendant_count|. That is
  iceberg selection under a different name, and the \fca{} literature has
  better selection criteria.
\item \textbf{Which concepts to materialise in a query planner.} The topic of
  Section~\ref{sec:dir-materialisation}.
\end{enumerate}

\paragraph{(c) What it would cost.} Stability is expensive to compute exactly
(it is a sum over subsets; the experiment script brute-forces it only for
extents of at most 16). Approximations exist and are standard practice. Since
the output is a ranking used for suggestion and indexing, approximation is
harmless --- it steers effort, never correctness.

\paragraph{(d) Verdict: \textbf{Adopt (derived)}.} Support is already there;
add estimated stability to the derived index, and use both to rank
suggestions. Never let either affect the stored graph.

\subsection{Boolean factor analysis and the MDL connection}
\label{sec:learn-mdl}

\paragraph{(a) What it is.} \citet{belohlavek2010} proved a striking result:
if you want to decompose a binary matrix into a product of two binary matrices
with the fewest inner dimensions --- Boolean factor analysis --- then the
optimal factors are \emph{formal concepts}. Concepts are not just an
interpretive device; they are the right basis.

That result immediately raises the model-selection question, because using
\emph{all} concepts is the trivial and useless decomposition. Which subset?
The minimum description length principle \citep{rissanen1978,grunwald2007}
gives a principled answer: keep a concept if describing the data using it is
shorter than describing the data without it --- if it \emph{pays for its own
description}. This is the programme of \citet{miettinen2014} and, in the
DAG-structured form relevant here, of the \texttt{mdl-fca} project
\citep{mdlfca}. The same two-part-MDL move drives the pattern-set miners
Krimp and Slim \citep{vreeken2011,smets2012}, which select a small set of
itemsets that compress the data; what a DAG-structured version adds is that
a new concept sits \emph{above} its constituents, and can make earlier ones
redundant.

\paragraph{(b) What \odag{} would gain.} A principled answer to the question
\odag{} currently punts on: \emph{where do categories come from?} Today they
come from people. MDL-selected concepts would let a store propose new
categories that are justified by compression --- Barlow's ``suspicious
coincidences'' \citep{barlow1989} made exact.

This is also the cleanest resolution of the tension in
Section~\ref{sec:not-closed}. The complaint against closure was that \fca{}
creates a concept for every coincidence. MDL is precisely a filter for
coincidence: a concept that arises from chance co-occurrence does not
compress, so it does not pay, so it is not admitted. \emph{MDL-\fca{} is the
principled version of iceberg partiality} --- support thresholds ask ``how
many?'', MDL asks ``does it earn its keep?''.

\paragraph{(c) What it would cost.} The learned structure is \emph{global and
history-dependent}: it is a function of the whole dataset, and it changes when
the data changes. Merging two learned DAGs would be merging two models fitted
to different samples, which is not union and does not converge. Putting
learned concepts directly into the shared graph would break G1 and G5 at once.

\paragraph{(d) Verdict: \textbf{Adopt (derived)}, with a clearly specified
promotion path.} The layering that makes it safe already exists:

\begin{itemize}[leftmargin=1.4em]
\item run the learner over a \emph{pinned root} --- so its input is
  immutable, its output is reproducible, and both can be cited;
\item keep the learned structure in a separate store with its own root, never
  merged (the same pattern as cone summaries);
\item if a learned category is to become shared knowledge, it enters as an
  \emph{ordinary asserted claim by an identified author}, with the pinned root
  it was derived from recorded as its basis. Someone --- a person or an
  accountable agent --- takes responsibility for it, and it becomes an
  assertion like any other.
\end{itemize}

Notice that this preserves the discipline of
Section~\ref{sec:learn-implications} exactly: the learner produces evidence,
an author produces claims.

\subsection{Attribute exploration, with a language model as the expert}
\label{sec:learn-exploration}

This is the most interesting item on the list, because a forty-year-old
protocol has just acquired a participant it was always waiting for.

\paragraph{(a) What it is.} Attribute exploration \citep{ganter2016} is an
interactive knowledge-acquisition loop:

\begin{enumerate}[leftmargin=1.7em,itemsep=1pt]
\item The system computes the next implication that its current data supports
  and that does not already follow from what has been confirmed.
\item It asks a domain expert: \emph{is this always true?}
\item If the expert says yes, the implication is added to the knowledge base.
\item If the expert says no, they must supply a \textbf{counterexample} --- an
  object violating it --- which is added to the context.
\item Repeat until no unconfirmed implications remain.
\end{enumerate}

On termination you have an implication base that is complete for the
\emph{domain}, not merely for the sample --- which is a strictly stronger
thing than data mining can give you, and the reason the protocol was invented.

\paragraph{(b) What \odag{} would gain.} A principled, terminating,
\emph{auditable} procedure for growing an ontology, in which the expensive
step --- answering domain questions --- is exactly the step a modern language
model is good at, and in which every answer is a discrete, recordable,
attributable commitment.

Consider how well the pieces fit:

\begin{center}\small
\begin{tabularx}{\textwidth}{@{}p{4.6cm} X@{}}
\toprule
\textbf{What exploration needs} & \textbf{What \odag{} has}\\
\midrule
a context to mine implications from & the asserted graph, viewed as
  $\ctx(P)$ (\S\ref{sec:dictionary})\\
an expert who answers questions & an agent on the MCP surface
  (\S\ref{sec:agents})\\
counterexamples added as objects & \cmd{put} --- the ordinary write path\\
a record of what was confirmed & signed provenance records, per claim\\
a way to revisit a bad answer & retraction records; roots to compare against\\
termination & the base is finite\\
\bottomrule
\end{tabularx}
\end{center}

And it addresses the central problem of agent-written knowledge, which
Section~\ref{sec:agents-review} states as: \emph{agents write faster than
anyone can review}. Exploration inverts the flow. Instead of an agent
proposing arbitrary edits for review, the \emph{system} generates the
questions --- a minimal, non-redundant sequence, each one maximally
informative given the answers so far --- and the agent answers them. Review
becomes bounded, ordered and non-redundant, which is exactly what makes
oversight tractable.

\paragraph{(c) What it would cost.} Three real problems, none fatal.

First, \textbf{the expert can be wrong}, and a language model can be
confidently wrong. Classical exploration assumes an infallible expert. The
mitigation is already \odag{}'s design: an answer is not truth but a
\emph{signed claim}, so a wrong answer is attributable, disputable and
retractable, and readers may weigh authors differently. The protocol degrades
from ``complete for the domain'' to ``complete relative to what this author
asserted'', which is the honest thing for it to be.

Second, \textbf{implications are not subsumptions}. Attribute exploration
yields $B \to C$ for sets $B$, and \odag{} stores only single-parent edges
$x \dsub y$. A confirmed implication with a compound premise
($\{\textsf{Flight}, \textsf{Japan}\} \to \textsf{Booked}$) has no
representation in the core graph. This is the relations/axioms wall, and it is
where the honest answer is: store the implication in the provenance layer as a
claim, use it as a lint, and do not pretend the core entails it.

Third, \textbf{exploration assumes a fixed universe}, while \odag{} stores
merge concurrently. Two agents exploring two replicas will ask different
questions and confirm different things. Reconciling that is genuinely open.

\paragraph{(d) Verdict: \textbf{Research}, and the most promising item in this
paper.} The single-writer case --- one store, one agent, one human reviewing
the confirmed implications --- is buildable today with existing pieces, and
would be a genuinely novel artefact: \emph{an ontology grown by a language
model under a protocol that bounds and orders the review}. The multi-writer
case needs the concurrent-exploration question answered first.
Section~\ref{sec:dir-exploration} sketches the design.

\subsection{Drawing: line diagrams as a design discipline}
\label{sec:learn-drawing}

\paragraph{(a) What it is.} \fca{} takes diagram layout seriously, because the
whole point of Wille's programme was that a domain expert should be able to
\emph{read} a lattice \citep{wille1982}. The techniques --- additive layouts
where each attribute contributes a fixed vector, reduced labelling
(\S\ref{sec:basic-theorem}), and layer assignment by concept rank --- are
mature.

\paragraph{(b) What \odag{} would gain.} Better pictures, and one specific
improvement worth naming. \odag{}'s visualiser draws the graph with Graphviz,
labelling every node with its full name. For a store with typed values that is
noisy: fifty weight values under one head, each drawn as a box. Reduced
labelling is the direct fix --- draw the structure, and place each name once,
where it is introduced. Compare Figure~\ref{fig:lattice} with
Figure~\ref{fig:ontodag}: the former carries the same information in less
space, and the difference is entirely the labelling convention.

\paragraph{(c) What it would cost.} Nothing; it is a rendering choice.

\paragraph{(d) Verdict: \textbf{Adopt (surface)}.} Cheap, and the reading
rules would have to be documented, which is a small cost against a real gain
in legibility.

\subsection{Algorithms and complexity as guardrails}
\label{sec:learn-complexity}

\paragraph{(a) What it is.} \fca{} has decades of algorithm engineering
\citep{kuznetsov2002,krajca2010,lindig2000} and a clear complexity picture
\citep{kuznetsov2001}.

\paragraph{(b) What \odag{} would gain.} Chiefly, \textbf{knowing where the
cliffs are}. The contranominal scale (\S\ref{sec:algorithms}) is a small,
memorable object whose lattice is $2^n$; recognising that pattern in a
proposal is a fast way to reject it. NextClosure's constant-memory,
lectic-order enumeration is the right algorithm for any future \odag{} feature
that must walk closed sets without materialising them --- an
implication lint, for instance, or an MDL learner's candidate generation.

There is a second, subtler gift. \fca{}'s complexity results explain
\emph{why} \odag{}'s design is not merely a convenient shortcut. A system that
promised ``all concepts'' would be promising something $\#\mathrm{P}$-hard to
even count. \odag{}'s refusal is principled, and the citation for that
principle is \citet{kuznetsov2001}.

\paragraph{(d) Verdict: \textbf{Adopt (theory)}.} Free.

\subsection{Triadic \fca{}, and provenance}

\paragraph{(a) What it is.} \citet{lehmann1995} extend contexts to three
dimensions: objects $\times$ attributes $\times$ \emph{conditions}, with
triadic concepts as triples.

\paragraph{(b) The tempting analogy.} \odag{}'s provenance layer records
\emph{who} asserted \emph{which} claim: a third axis over the same claims.
``Claims that Alice and Bob both endorse'' looks like a triadic concept.

\paragraph{(c) Why it does not fit as-is.} Triadic \fca{} treats the third
dimension symmetrically with the other two, and \odag{}'s does not:
authorship is not a property of knowledge but of a speech act, deliberately
stored in a \emph{separate} store with its own root so that the knowledge root
stays a fingerprint of knowledge alone. Merging the axes would put ``who said
it'' inside the hash of ``what is true'', which is precisely the design error
the separation exists to prevent --- two people who know the same thing must
get the same root.

\paragraph{(d) Verdict: \textbf{Adopt (theory), refuse (core)}.} The triadic
\emph{analysis} could run over the provenance store as a derived view --- ``which
authors agree about which regions of the ontology'' is a real and interesting
question, and triadic concepts are the right shape for it. The triadic
\emph{representation} must stay out of the knowledge graph.

\subsection{What \odag{} should refuse}

Three items, for completeness.

\paragraph{Fuzzy and graded \fca{}} \citep{belohlavek2004}. \textbf{Refuse.} A
graded incidence puts a number in the data. Numbers in identity break exact
canonicalisation (which number, at what precision, rounded how?) and break
merge (whose number wins?). Confidence is a property of a \emph{speech act},
and belongs on a provenance record where it can be attributed and where
disagreement is representable as disagreement rather than as a silently
averaged scalar.

\paragraph{Full closure in the shared store.} \textbf{Refuse}, for the four
reasons of Section~\ref{sec:not-closed}. Local closure over a pinned root is
fine and is the recommended way to get the benefits.

\paragraph{Relational Concept Analysis as the relations answer}
\citep{rouanehacene2013}. \textbf{Research.} \odag{} genuinely lacks
relations: ``Alice \emph{authored} Doc1'' is not subsumption, and the system
has no way to say it. RCA is the mature \fca{}-native answer, handling
inter-object links by iterated relational scaling to a fixpoint. But the
fixpoint is a global computation over the whole context, which is the opposite
of \odag{}'s local, incremental, mergeable discipline --- and the scaled
attributes it generates would enter the canonical form, making the root depend
on the fixpoint. Worth reading before designing anything here; not worth
importing.

\subsection{The shopping list, ranked}

\begin{table}[h]\centering\small
\caption{Everything in Section~\ref{sec:learn}, sorted by value per unit of
risk. The top three cost nothing and can be done now.}
\label{tab:shopping}
\vspace{1mm}
\begin{tabularx}{\textwidth}{@{}p{4.4cm} p{2.5cm} X@{}}
\toprule
\textbf{Item} & \textbf{Verdict} & \textbf{Why, in one line}\\
\midrule
Pattern structures (\S\ref{sec:learn-patterns})
  & Adopt (theory) & \odag{} already built one; the theory says what is true
    about it and how to extend it to several coordinates.\\
Complexity results (\S\ref{sec:learn-complexity})
  & Adopt (theory) & Turns a design instinct into a citable argument, and
    names the cliffs.\\
Scaling vocabulary (\S\ref{sec:learn-scaling})
  & Adopt (core) & Reveals a real gap: there is no ordinal kind.\\
Reduced labelling (\S\ref{sec:learn-drawing})
  & Adopt (surface) & Strictly better pictures for free.\\
Stability \& support (\S\ref{sec:learn-stability})
  & Adopt (derived) & Half of it already exists; ranks which meets and cones
    deserve attention.\\
Implication lint (\S\ref{sec:learn-implications})
  & Adopt (derived) & Real classification help; must stay a suggestion,
    because implications are non-monotone.\\
MDL concept learning (\S\ref{sec:learn-mdl})
  & Adopt (derived) & Answers ``where do categories come from?''; needs a
    pinned root and an author to promote its output.\\
Attribute exploration (\S\ref{sec:learn-exploration})
  & Research & The best fit with the agents story; single-writer case is
    buildable now.\\
Triadic analysis of provenance
  & Research & Right shape for ``who agrees with whom''; must not touch the
    knowledge root.\\
Relational Concept Analysis
  & Research & The relations wall; global fixpoint conflicts with local
    canonical form.\\
Fuzzy \fca{}
  & Refuse & Numbers in identity destroy exact canonicalisation and merge.\\
Full closure in the store
  & Refuse & Exponential, coincidence-prone, unstable under edit,
    unattributable.\\
\bottomrule
\end{tabularx}
\end{table}

The pattern across the table is worth stating explicitly, because it is the
answer to ``how much of \fca{} should \odag{} adopt?''.

\begin{keyidea}
\textbf{Almost all of \fca{}'s theory is adoptable; almost none of its
representation is.} The theorems, the vocabulary, the algorithms, the
selection criteria and the interactive protocol can all be imported. The one
thing that cannot is the concept lattice itself as a stored object --- because
storing it would forfeit the cheap canonical form, and the cheap canonical
form is the thing everything in Part~IV is built on.
\end{keyidea}
